% Qwen3.5 на Android: GPU и NPU.
% Сборка: latexmk main.tex (движок lualatex, см. .latexmkrc)
% Экспорт в docx: python3 tex2docx.py
\documentclass[11pt]{article}

\usepackage[a4paper,top=2.3cm,bottom=2.3cm,left=2.4cm,right=2.4cm]{geometry}
\usepackage{qwen}
\usepackage[colorlinks=true,linkcolor=artAccent,urlcolor=artAccent]{hyperref}
% Ссылка на репозиторий длиннее остатка строки: xurl разрешает перенос внутри
% адреса, иначе строка выходит за поле.
\usepackage{xurl}

\title{\vspace{-1.2em}\displayfamily\LARGE\bfseries
Qwen3.5 на Android: GPU и NPU\vspace{0.3em}}
\author{Григорий Евко\\[0.25em]
\normalsize\mdseries\url{https://github.com/GrigoryEvko/qwen-mobile}}
\date{}

\begin{document}
\maketitle
\vspace{-1.8em}

% --------------------------------------------------------------- 1
\section{Результаты}

Измерения выполнены на OnePlus~13 на холодном телефоне (тепловом статусе~0,
без кабеля, утилитой \texttt{llama-bench}): промпт 512 токенов, ответ 64
токена, три повтора. Спекулятивное декодирование измерено отдельно, на
свободном промпте в 96 токенов.

\begin{qtable}{l l N{4}{1} N{2}{1} N{2}{1}}{Скорость обработки промпта и
генерации ответа, токенов в секунду. Последний столбец --- генерация со
спекулятивным декодированием.}{tab:speed}
\thead{Модель} & \thead{Где} & {\thead{Промпт}} & {\thead{Ответ}} & {\thead{Спекуляция}} \\ \midrule
\tstub{2B, F16} & CPU, 8 потоков & 67.6 & 12.9 & {--} \\
\tstub{2B, F16} & GPU Adreno & 271.6 & 13.6 & {--} \\
\tstub{2B, F16} & NPU Hexagon & 1456.0 & 13.8 & {--} \\ \addlinespace
\tstub{2B, Q8\_0} & NPU Hexagon & 1558.0 & 23.6 & 28.8 \\
\tstub{4B, Q8\_0} & NPU Hexagon & 755.0 & 10.4 & 15.4 \\ \addlinespace
\tstub{2B, Q4\_0} & NPU Hexagon & 1022.0 & 39.1 & {--} \\
\end{qtable}

Изображение на 576 токенов кодируется на NPU (нейропроцессор, он же Hexagon
DSP) за 0{,}55--0{,}82~с.

Скорость декода упирается в пропускную способность памяти, измеренная скорость
чтения составляет 45--51~ГБ/с, у 4B пары матриц \texttt{gate} и \texttt{up}
читаются на 56~ГБ/с. Префилл работает относительно хорошо: 512 токенов 2B
требуют 1{,}41~TFLOP, что при 1558 токенах в секунду даёт около 4{,}4~TFLOPS.

Качество измерялось как KL-дивергенция на токен относительно F16-логитов
исходного чекпойнта, вычисленных на CUDA (WikiText-2, 16 отрезков по 512
токенов). Продуктовые файлы собраны собственным упаковщиком из исходных
чекпойнтов: равномерный Q8\_0 без вращений, свёрток и калибровки.

\begin{qtable}{l N{1}{2} N{1}{5} N{1}{4} N{1}{2} N{2}{2}}{Расхождение с F16 и
пол точности каждого бэкенда.}{tab:kl}
\thead{Файл или бэкенд} & {\thead{ГиБ}} & {\thead{Средняя KL}} & {\thead{99{,}9\,\%}} & {\thead{Макс.}} & {\thead{Топ-1, \%}} \\ \midrule
\tstub{2B Q8\_0} & 1.99 & 0.00115 & 0.0143 & 0.03 & 98.21 \\
\tstub{4B Q8\_0} & 4.41 & 0.00202 & 0.1750 & 1.56 & 97.21 \\ \addlinespace
\tstub{F16 на CPU} & {--} & 0.00020 & 0.0016 & 0.00 & 99.09 \\
\tstub{F16 на NPU} & {--} & 0.00019 & 0.0015 & 0.00 & 99.12 \\
\tstub{F16 на GPU} & {--} & 0.00065 & 0.0078 & 0.01 & 98.53 \\
\end{qtable}

NPU численно неотличим от CPU, переход на GPU обходится дороже, чем
квантизация в 8 бит, скорее всего там баги в бэкенде.

% --------------------------------------------------------------- 2
\section{Баги}

\subsection{Чтение bias по невыровненному адресу}

В бэкенде Hexagon есть несколько десятков флагов переключения режима
исполнения. При включённом флаге слияния операций бэкенд читал прибавляемый
вектор с неверно выровненными значениями. Строки матрицы делятся между
потоками чётным числом, а не кратным~32, поэтому смещение среза не кратно
128~байтам, и
все потоки, кроме первого, читали чужие данные. В Qwen3.5 на это попадал
вектор \texttt{ssm\_dt.bias}, корраптился гейт затухания рекуррентного слоя, и
текст расходился начиная со второго токена.

\subsection{Неверная форма GELU}

Операции с фронтенда маршрутизировались на сигмоидный GELU из бэкенда,
кодировщик изображений Qwen3.5 использует точную форму через tanh в 24 слоях.

% --------------------------------------------------------------- 3
\section{Оптимизации}

На батче из 512 токенов операция gated delta net занимала 235~мс из 501, то
есть 44\,\% всей обработки промпта, и шла со скоростью 0{,}12~TFLOPS.
Рекуррентность вычислялась токен за токеном, по 11 векторных операций на
вектор состояния.

Оптимизация: последовательность режется на блоки по 64 токена; внутри блока
рекуррентность сводится к произведениям матриц, которые выполняются на HMX
(матричный блок) в fp16 с накоплением в f32, а треугольная система решается
прямой подстановкой на HVX (векторный блок) в f32. Состояние остаётся в f32 в
VTCM (быстрая локальная память). Раскладка
занимает 450~КБ на одну строку данных, и все 16 голов слоя 2B помещаются в
8~МБ VTCM одной группой. В f32 расхождение с эталоном составляет $10^{-13}$, а
с округлением fp16 на границах HMX --- около $3\cdot10^{-7}$ и не растёт с
длиной последовательности, поскольку каждый блок умножает перенесённую ошибку
на множитель не больше единицы. На телефоне файл 2B Q8\_0 даёт максимальную KL
0{,}00986 с новым ядром против 0{,}00983 со старым. Обработка промпта 2B
выросла с 981 до 1558 токенов в секунду, 4B --- с 382 до 755.

Параллельно были устранены две другие потери на тех же слоях. При генерации
одного токена слой создавал около 14 мелких операций на каждое рекуррентное
состояние, 790 операций на батч; после оптимизации батч на DSP (он же NPU, включает HMX и HVX) сократился с
31{,}0 до 24{,}8~мс на токен при 513 операциях вместо 790. На батче пара
операций конкатенации и свёртки переносила токены дважды: первая писала копию
всего батча, вторая её читала.

\subsection{Оверхед запуска и спекулятивное декодирование}

Сам DSP считает токен 24{,}8~мс, а весь токен занимает около 29~мс --- значит,
4{,}4--4{,}7~мс приходится на хост. Передача задания в DSP и ожидание ответа
из них занимают только 0{,}3~мс, сборка графа в llama.cpp --- ноль, на одном
токене граф переиспользуется. Остальное хост тратил так: бэкенд заново разбирал
все 513 операций на каждом вызове, планировщик делал лишний проход на CPU ради
выборки строки из таблицы эмбеддингов, сэмплер перебирал все
кандидаты, чтобы оставить двадцать.

Сессия Hexagon создаётся на устройство, а не на контекст, и хранила разобранный граф
в единственной ячейке. Драфт-модель и основная --- это два разных графа на
одной сессии, поэтому каждый шаг спекуляции вытеснял предыдущий и заново
собирал 513 дескрипторов. Спекуляция из-за него проигрывала: 17{,}3 токена в секунду
против 21{,}1 без неё. После замены ячейки на кеш по идентификатору графа и
добавления снимка упакованных дескрипторов спекуляция стала выигрывать на
обеих моделях: 2B --- 28{,}8 против 23{,}6, 4B --- 15{,}4 против 10{,}4.

Выигрыш зависит от текста: на перечислениях и шаблонных ответах драфт
принимается в 88--92\,\% случаев, на свободном --- в 40--71\,\%. Шаг с драфтом
в три токена стоит 2{,}1 обычного шага и даёт $1+3a$ токенов, где $a$ --- доля
принятых, то есть окупается начиная с $a\approx0{,}37$; в приложении длина
драфта подстраивается под наблюдаемую долю по окну из 32 последних шагов.


% --------------------------------------------------------------- 4
\section{Квантизация}

У 4B в Q8\_0 нашёлся тяжёлый хвост: восемь токенов из 8192 дают KL выше
0{,}17 при максимуме 1{,}56, тогда как у 2B максимум 0{,}028. Перебором найнден \texttt{ssm\_out}, выходная проекция
рекуррентных слоёв. Если оставить только её в F16, максимум падает до
0{,}553, а квантиль 99{,}9\,\% --- с 0{,}175 до 0{,}037 ценой 225~МиБ. На
мегабайт добавленного размера этот класс убирает хвост в 3{,}5 раза
эффективнее таблицы эмбеддингов и в 14 раз эффективнее всего MLP. Тот же класс
оказался первым и на 2B в четырёхбитных экспериментах, и в обзоре unsloth
модели 35B-A3B.

Четырёхбитная часть пайплайна (свёртка норм, вращение остаточного потока
матрицей Адамара, перестановка каналов MLP, решатель GPTQ или Qronos, блочная
оптимизация) доведена на 2B до 1{,}62~ГиБ при средней KL 0{,}0295 и согласии
топ-1 91{,}4\,\%. Лестница экспериментов показала, что округление обязано
происходить решателем по Гессиану, а веса при блочной оптимизации должны быть
заморожены, иначе результат переобучается на калибровочной выборке.
Дополнительно голова модели привязывается обратно к таблице эмбеддингов одной
плотной матрицей, что даёт файл в 1{,}13~ГиБ при KL 0{,}0311. Этот путь пока в
резерве, продукт собран на Q8\_0.

Ниже четырёх бит спуститься нельзя: ни один трёхбитный формат ggml не
исполняется ни на Hexagon, ни на Adreno, такой файл работал бы только на CPU.

% --------------------------------------------------------------- 5
\section{Приложение и сборка}

Сделал для тестов приложение qwen-mobile
(\url{https://github.com/GrigoryEvko/qwen-mobile}, там же релизы): чат с
парсингом Markdown и сворачиваемым блоком
рассуждений, работа с изображениями, настройки (выбор ускорителя, длина
контекста, детализация изображений, спекуляция), вкладка отладки с бенчмарком
и измерением энергии, сохранение диалогов и кеш состояний, за счёт которого
повторная отправка того же изображения и продолжение диалога не требуют
повторных вычислений. Выбор ускорителя --- CPU, GPU, NPU. Модели не
поставляются и не скачиваются: файл GGUF и проектор изображений кладутся в
каталог на устройстве.

% --------------------------------------------------------------- 6
\section{Таски}

\begin{itemize}
\item \textbf{Совместимость спекуляции со слияниями.} Описана выше, поднимает
и префилл, и генерацию в режиме с драфтом.
\item \textbf{Кодировщик изображений.} Внимание при размерности головы 64 идёт
на трети той скорости, которую чип выдаёт на умножении матриц, и занимает
38\,\% времени кодирования. Ожидаемый выигрыш --- примерно с 0{,}7 до
0{,}45~с на изображение.
\item \textbf{Энергия на токен.} Измерение встроено в приложение, таблицы по
ускорителям пока нет. Это же ответ на вопрос про нагрев: в обычном диалоге
телефон не троттлит и держит тепловой статус~0, 70\,°C были получены один раз
на десятиминутном непрерывном счёте с подключённым кабелем.
\item \textbf{Другие чипсеты.} Библиотека для DSP собирается из того же дерева
для версий v73, v75, v79 и v81, выбор происходит во время выполнения, но в
текущий пакет вложена только сборка под v79. Для 8~Gen~3 нужен пересбор с
одной изменённой строкой.
\item \textbf{Перемеры.} Базовые числа F16 для CPU и GPU и скорость 39{,}1
токена в секунду для 2B Q4\_0 взяты из рабочих записей и должны быть получены
заново \remeasure{в журнале измерений репозитория их нет}. Сравнение времени
кодирования изображения на GPU и NPU сделано при разном числе токенов
\remeasure{7~с на GPU измерены при 672 токенах, 0{,}55--0{,}82~с на NPU --- при
576}.
\end{itemize}

\end{document}
